%===============================================================================
% 双向同步与编译说明
% 1. 修改“星尘黑客松-中文简历.tex”后，请同步更新同目录下的“星尘黑客松-中文简历.md”。
% 2. AI 编辑本文件时，如发现 .tex 与 .md 内容不一致，应立即提醒用户并说明差异。
% 3. 需要生成 PDF 时，请使用 TeXPage、Overleaf 等在线 LaTeX 工具，并将编译器设置为 XeLaTeX。
% 4. 用途：星尘智能 Astribot OS「一句话生成App」机器人黑客松报名/展示用中文简历（与正式学术向「中文简历」成对文件不同）。
%===============================================================================

%-------------------------
% 李明哲 - 星尘黑客松中文简历
% XeLaTeX 专用版本
%------------------------

\documentclass[letterpaper,12pt]{article}

\usepackage{latexsym}
\usepackage[empty]{fullpage}
\usepackage{titlesec}
\usepackage{marvosym}
\usepackage[usenames,dvipsnames]{color}
\usepackage{verbatim}
\usepackage{enumitem}
\usepackage[hidelinks]{hyperref}
\usepackage{fancyhdr}
\usepackage{tabularx}

%----------中文字体设置 (XeLaTeX only)----------
\usepackage{fontspec}
\usepackage{xeCJK}
\setCJKmainfont{SimSun} % 宋体
\setmainfont{Arial}

%----------FONT OPTIONS----------
\usepackage[default]{sourcesanspro}

\pagestyle{fancy}
\fancyhf{}
\fancyfoot{}
\renewcommand{\headrulewidth}{0pt}
\renewcommand{\footrulewidth}{0pt}

% Adjust margins
\addtolength{\oddsidemargin}{-0.5in}
\addtolength{\evensidemargin}{-0.5in}
\addtolength{\textwidth}{1in}
\addtolength{\topmargin}{-.55in}
\addtolength{\textheight}{1.1in}

\urlstyle{same}

\raggedbottom
\raggedright
\setlength{\tabcolsep}{0in}

% Sections formatting
\titleformat{\section}{
  \vspace{-4pt}\scshape\raggedright\large
}{}{0em}{}[\color{black}\titlerule \vspace{-5pt}]

%-------------------------
% Custom commands
\newcommand{\resumeItem}[1]{
  \item\small{
    {#1 \vspace{-2pt}}
  }
}

\newcommand{\resumeSubheading}[4]{
  \vspace{-2pt}\item
    \begin{tabular*}{0.97\textwidth}[t]{l@{\extracolsep{\fill}}r}
      \textbf{#1} & #2 \\
      \textit{\small#3} & \textit{\small #4} \\
    \end{tabular*}\vspace{-7pt}
}

\newcommand{\resumeProjectHeading}[2]{
    \item
    \begin{tabular*}{0.97\textwidth}{l@{\extracolsep{\fill}}r}
      \small#1 & #2 \\
    \end{tabular*}\vspace{-7pt}
}

\renewcommand\labelitemii{$\vcenter{\hbox{\tiny$\bullet$}}$}

\newcommand{\resumeSubHeadingListStart}{\begin{itemize}[leftmargin=0.15in, label={}]}
\newcommand{\resumeSubHeadingListEnd}{\end{itemize}}
\newcommand{\resumeItemListStart}{\begin{itemize}[leftmargin=0.22in]}
\newcommand{\resumeItemListEnd}{\end{itemize}\vspace{-5pt}}

%-------------------------------------------
%%%%%%  RESUME STARTS HERE  %%%%%%%%%%%%%%%%%%%%%%%%%%%%

\begin{document}

%========================
% 中文简历
%========================

%----------HEADING----------
\begin{center}
    \textbf{\Huge \scshape 李明哲} \\ \vspace{5pt}
    \small 机器人系统 / 具身智能 / 场景App与任务智能体 \\ \vspace{3pt}
    \small 邮箱：mc64655 [at] um [dot] edu [dot] mo $|$ 手机/微信：(+86) 137 1041 1101 $|$ 当前所在地：深圳 \\
    \small 个人主页：\href{https://mzli112358.github.io/}{https://mzli112358.github.io/}
\end{center}

%-----------SUMMARY-----------
\section{个人简介}
 \begin{itemize}[leftmargin=0.15in, label={}]
    \small{\item{
    面向具身智能、家庭服务机器人与真实场景交付的全栈机器人开发者，具备从方案设计、硬件原型、感知导航、遥操作数据采集到VLA适配与任务智能体编排的端到端经验。当前在原力无限科技控股（浙江）有限公司深圳研究院参与家庭服务机器人研发。曾创立35人高校机器人实验室并获RoboMaster相关奖项。个人报名星尘Astribot OS黑客松，希望把此前因平台执行权限受限而未能完成的PhotoMate方案，在星尘本体与OS能力生态上真正做成可运行、可演示、可复用的具身智能摄影师App。
    }}
 \end{itemize}

%-----------CURRENT EXPERIENCE-----------
\section{近期经历}
    \resumeSubHeadingListStart

    \resumeSubheading
      {原力无限科技控股（浙江）有限公司 / INFIFORCE 深圳研究院}{深圳，中国}
      {科研助理 / 具身智能系统工程实习生}{2026年3月 -- 至今}
      \resumeItemListStart
        \resumeItem{参与家庭服务机器人与VLA系统研发，围绕“机器人本体部署—真实数据采集—模型训练—任务执行反馈”闭环，推进硬件原型、遥操作、数据采集与模型适配工作。}
        \resumeItem{负责从0到1快速搭建机器人原型：完成方案调研、展会与供应链调研、机械结构设计、零部件选型、电路设计与打板、焊接装配、线束整理、驱动适配、电机校零、CAN通信与整机联调。}
        \resumeItem{参与多种机器人本体与机械臂平台开发和适配，包括松灵Piper、Nero、OpenArm、Franka、星海图R1等，熟悉机器人硬件调试、ROS通信、控制接口适配与数据采集流程。}
        \resumeItem{搭建并缝合SLAM、定位、建图、导航、视觉感知、相机标定、VR遥操作、数据采集与模型训练相关开源仓库，推动机器人从单臂、双臂、夹爪换手到全身系统的能力扩展。}
        \resumeItem{熟悉LeRobot框架，参与ACT模型微调与OpenPI $\pi_{0.5}$模型微调，用于机器人操作任务的数据采集、策略学习与跨本体适配。}
        \resumeItem{参与长程任务智能体系统建设，负责智能体任务管理、自主导航定位、任务规划与执行链路设计，协调两名智能体方向成员推进开发。}
        \resumeItem{承担小团队工程管理与技术审查职责，组织机械、电控、VR遥操作、智能体等方向实习生成员协作，进行方案评审、图纸评审、PCB评审、任务拆解、进度推进与联调问题定位。}
        \resumeItem{具备高强度快速响应与持续迭代能力，能够在硬件、软件、算法与产品需求之间快速取舍，推动原型从概念验证走向可演示系统。}
      \resumeItemListEnd

    \resumeSubHeadingListEnd


%-----------ASTRIBOT HACKATHON FIT-----------
\section{星尘黑客松相关能力}
    \resumeSubHeadingListStart
        \resumeItem{体验驱动交付：从“用户想获得什么体验”出发定义服务闭环，契合Astribot OS由自然语言需求组合机器人能力、生成场景App的路线。}
        \resumeItem{场景App编排：具备任务智能体、状态机与工具调用经验，可将导航、感知、多模交互、灵巧操作、拍摄及交付模块编排为可恢复的长流程。}
        \resumeItem{本体无关设计：PhotoMate按“能力接口—任务状态—外设适配”分层，不绑定单一机器人；更换本体时主要重做运动能力映射、坐标标定与外设适配。}
        \resumeItem{真机快速落地：能够在短周期完成需求拆解、机械结构、3D打印、外设集成、软件缝合、整机排障与演示设计。}
    \resumeSubHeadingListEnd

%-----------RELEVANT PROJECT-----------
\section{相关项目}
    \resumeSubHeadingListStart
    \resumeSubheading
      {PhotoMate：面向商业空间的具身智能摄影师}{2026年7月}
      {项目负责人 / 系统设计与集成}{探月黑客松方案延续，拟基于Astribot OS重启}
      \resumeItemListStart
        \resumeItem{\textbf{问题与场景：}面向展会、景区、毕业典礼和商场，解决无人帮拍、自拍构图受限，以及传统摄影成本高、覆盖不足的问题；机器人不是固定自拍亭，而是能主动发现需求、移动服务并创造互动体验的摄影师。}
        \resumeItem{\textbf{完整流程：}巡航/待命—感知驻足、招手或语音请求并主动邀请—选择用户手机/机器人相机及打印/录像—导航至预设机位—分析人体、人脸、背景和光线并调整相机姿态—语音指导与倒计时—拍摄质检及失败重试—二维码或实体照片交付—清理数据并返回巡航。}
        \resumeItem{\textbf{两种模式：}用户手机模式以末端夹具稳持设备、灵巧手触屏快门，照片不离开用户手机；机器人相机模式接入Insta360等USB相机，支持构图、曝光、闭眼/模糊检测、选片、短视频和扫码交付。}
        \resumeItem{\textbf{实现架构：}按本体与外设、感知、能力、App编排四层设计；封装导航、机械臂/灵巧手、相机、语音、显示和打印工具，由显式状态机/任务智能体管理分支、超时、重试、恢复和资源释放，而非让大模型直接控制全部硬件。}
        \resumeItem{\textbf{跨本体设计：}核心App仅依赖移动到机位、调整相机位姿、语音交互、触发拍摄、操作/递交等标准能力接口；不同机器人通过适配层完成能力映射、坐标标定与外设接入。}
        \resumeItem{\textbf{此前受阻：}探月社区Physical AI黑客松期间，赞助商本体未开放导航、机械臂/灵巧手等所需执行权限，无法完成关键动作调用与真机闭环；项目停留在产品方案、结构/外设设计和部分软件链路，未形成完整Demo。}
        \resumeItem{\textbf{星尘重启：}Astribot OS的模型—系统—本体一体化架构、Meta能力包、SDK/ROS2接口与应用编排思路，能够补足此前的本体调用断点；计划借助星尘真机、OS和开发支持，将PhotoMate实现为可运行App及可复制的摄影半解决方案。}
      \resumeItemListEnd
    \resumeSubHeadingListEnd


%-----------ROBOTICS EXPERIENCE-----------
\section{机器人与工程经历}
    \resumeSubHeadingListStart

    \resumeSubheading
      {北师香港浸会大学 Navigator Robotics Lab}{珠海，中国}
      {创始人 / 技术负责人}{2023年4月 -- 至今}
      \resumeItemListStart
        \resumeItem{创立学校首个机器人实验室，推动团队从0到1建设，获得约20万元人民币经费支持，团队规模扩展至35人。}
        \resumeItem{负责机器人系统全栈集成，包括工业相机、激光雷达、Jetson计算平台、ROS2通信与实时控制系统，用于对抗机器人、轮式机械臂机器人等平台。}
        \resumeItem{带队获得RoboMaster机甲大师高校联盟赛工程挑战赛辽宁站第一名。}
        \resumeItem{具备机械设计、SolidWorks建模、3D打印、传感器融合、快速原型开发、现场调试与工程交付经验。}
      \resumeItemListEnd

    \resumeSubHeadingListEnd


%-----------INTERNSHIP EXPERIENCE-----------
\section{其他实习经历}
    \resumeSubHeadingListStart

    \resumeSubheading
      {深圳市威世博知识产权代理事务所 / 小美知识产权集团}{深圳，中国}
      {大语言模型工程实习生}{2025年6月 -- 2025年8月}
      \resumeItemListStart
        \resumeItem{构建基于LangChain的自动化多智能体系统，用于专利文档处理、流程自动化、知识检索与任务协同。}
        \resumeItem{参与GPU集群管理与模型训练流程，支持大模型实验、部署与工程化验证。}
      \resumeItemListEnd

    \resumeSubheading
      {艾普工华科技（武汉）有限公司 / EpicHust}{武汉，中国}
      {大语言模型工程实习生}{2025年1月 -- 2025年2月}
      \resumeItemListStart
        \resumeItem{参与面向工业协议与智能工厂场景的开源模型微调与评估，支持企业级AI应用落地。}
        \resumeItem{参与RAG问答、知识库构建与模型部署方案设计，面向工业文档、设备协议与企业知识检索场景。}
      \resumeItemListEnd

    \resumeSubHeadingListEnd

%-----------EDUCATION-----------
\section{教育经历}
  \resumeSubHeadingListStart
    \resumeSubheading
      {澳门大学}{澳门，中国}
      {机器人与自主系统理学硕士（拟入学）}{2026年8月 -- 2028年6月}
    \resumeSubheading
      {北师香港浸会大学}{珠海，中国}
      {数据科学理学学士（荣誉）}{2022年8月 -- 2026年6月}
  \resumeSubHeadingListEnd

%-----------SKILLS-----------
\section{专业技能}
 \begin{itemize}[leftmargin=0.15in, label={}]
    \small{\item{
     \textbf{具身智能与VLA}{: LeRobot、ACT策略微调、OpenPI $\pi_{0.5}$微调、机器人数据采集、遥操作、任务规划、长程任务智能体、多智能体系统} \\
     \textbf{机器人系统}{: ROS1/ROS2、Gazebo、SLAM、定位建图导航、CAN通信、电机驱动调试、机械臂控制、夹爪/末端执行器适配、传感器融合} \\
     \textbf{硬件原型}{: 机械结构设计、SolidWorks、OnShape、3D打印、PCB设计与打板、焊接装配、线束整理、整机联调、供应链调研与选型} \\
     \textbf{3D视觉与仿真}{: 3D Gaussian Splatting、NeRF、几何一致SLAM、NVIDIA Isaac Lab、机器人强化学习、物理仿真、相机标定、视觉感知} \\
     \textbf{机器学习与大模型}{: PyTorch、强化学习、Transformer、RAG、LangChain、模型微调、上下文学习、多智能体工作流} \\
     \textbf{编程与工具}{: Python、C++、CUDA、Linux、Git、Docker、FastAPI/Flask、前后端快速开发、CAD/FEA、3D打印} \\
     \textbf{使用/适配平台}{: 松灵Piper、Nero、OpenArm、Franka、星海图R1、Jetson平台、深度相机、激光雷达、VR遥操作设备} \\
     \textbf{语言能力}{: 英语，可用于学术阅读、技术文档阅读与专业沟通} \\
    }}
 \end{itemize}



\end{document}
